\section{Conclusion}
\label{sec:conclusion_future_work}

This thesis set out to investigate whether a lightweight neuro-symbolic anomaly-analysis framework could provide a more interpretable, operationally meaningful, and deployment-conscious alternative to purely statistical anomaly detection for weather-dependent electricity systems. The proposed solution combined a compact forecasting backbone based on dCeNN--ELM, residual-based anomaly generation, ASP-based symbolic refinement, and finance-aware downstream evaluation. Rather than treating anomaly detection as a purely statistical ranking problem, the research framed it as a layered decision pipeline in which learning, reasoning, and practical interpretation each play a distinct role.

Taken as a whole, the work supports a measured but positive conclusion. The hybrid approach does not uniformly dominate all baselines in pure forecasting accuracy, and therefore it should not be presented as a universally superior standalone predictor. However, that was never the strongest or most distinctive contribution of the thesis. The more important result is that the proposed framework demonstrates how a compact forecasting model, when coupled with explicit symbolic reasoning and finance-aware interpretation, can produce anomaly signals that are more interpretable, more plausibility-aware, and more suitable for downstream use than raw statistical detections alone. In that sense, the thesis successfully shows the value of combining learning and reasoning in a resource-conscious anomaly-analysis pipeline \cite{repo_readme_framework,repo_benchmark_comparison,repo_veto_analysis_plot,repo_finance_mapping}.

\subsection{Summary of Findings}
\label{subsec:summary_of_findings}

Several findings emerge clearly from the results.

First, the proposed framework provides a coherent \emph{end-to-end neuro-symbolic pipeline}. The repository structure and execution workflow demonstrate that the project is not merely a forecasting notebook or an isolated anomaly script, but a full pipeline from data engineering through model training, threshold calibration, anomaly detection, symbolic refinement, utility mapping, event construction, and edge export \cite{repo_readme_framework}. This is itself a meaningful contribution, because many anomaly-detection studies stop at raw prediction or point-level anomaly scoring.

Second, the forecasting analysis showed that the proposed dCeNN--ELM backbone should be interpreted as a lightweight nonlinear forecasting core rather than as a claim of universal predictive dominance. In the committed benchmark outputs, Ridge Regression is the strongest model in pure forecasting accuracy, while dCeNN--ELM remains substantially more lightweight than the LSTM baseline in training time, inference latency, and model size \cite{repo_benchmark_comparison}. This result is important because it locates the real contribution of the proposed neural design: compactness, modularity, and compatibility with later symbolic and edge-oriented stages.

Third, the symbolic layer contributes one of the strongest outcomes of the thesis: \emph{interpretable refinement}. The ASP stage does not simply reduce anomaly counts. It introduces explicit veto logic grounded in persistence, physical plausibility, and cross-signal support. This gives the system a principled way to reject weak or implausible anomaly candidates and therefore addresses one of the central shortcomings of purely residual-based detectors: the absence of a domain-grounded explanation for why an anomaly should or should not be trusted \cite{repo_veto_analysis_plot,repo_report_tables}.

Fourth, the financial-impact analysis suggests that anomaly refinement should be judged not only by how many alarms are generated, but by whether the retained anomaly signals are more useful downstream. The finance-aware mapping stage and neural-versus-neuro-symbolic utility comparison move the evaluation away from abstract anomaly counts and toward a more practical criterion of usefulness \cite{repo_finance_mapping,repo_utility_comparison}. Even when the forecasting layer is not numerically dominant, the full framework can still add value by improving the quality and usability of the signals that survive refinement.

Fifth, the event-based analysis demonstrates that the framework is capable of surfacing structured high-severity episodes rather than only isolated anomaly points. The event table, master timelines, and event comparisons support the view that the system can capture prolonged price disturbances, severe load deviations, renewable-side stress episodes, and multi-signal events in a coherent way \cite{repo_anomaly_events_2022,repo_build_event_table,repo_plot_master_timeline}. This strengthens the practical relevance of the thesis because real operators and analysts reason in terms of events, not only pointwise scores.

Taken together, these findings show that the hybrid approach is successful in the sense most relevant to the thesis objectives. Its success lies less in absolute forecasting dominance and more in its ability to integrate lightweight learning, symbolic rejection logic, event abstraction, and downstream utility interpretation into one coherent system.

\subsection{Limitations}
\label{subsec:limitations}

Although the thesis demonstrates the practical value of the hybrid approach, several limitations should be acknowledged clearly.

The first and most important limitation is the \emph{dependence on rule engineering}. The ASP refinement layer is only as strong as the domain knowledge encoded within it. While the current rule base captures persistence, solar plausibility, wind-ramp filtering, calendar context, and selected market logic, it remains intentionally compact. This means that the interpretability and false-alarm reduction achieved by the system are bounded by the completeness and quality of the implemented rules. In practical terms, the ASP layer does not automatically discover domain knowledge; it depends on careful human design. This creates a knowledge-engineering bottleneck and means that transferring the framework to another domain or market would require rule adaptation rather than direct plug-and-play reuse.

A second limitation concerns the forecasting layer. The results do not support a strong claim that dCeNN--ELM is the best forecaster among all tested alternatives. In the committed benchmark artifacts, Ridge performs more strongly in raw RMSE terms, and the detailed monthly metrics show that some targets---especially price and wind during volatile periods---remain difficult for the implemented forecasting backbone \cite{repo_benchmark_comparison,repo_metrics_monthly}. This implies that the anomaly-generation stage inherits a certain amount of forecasting instability, which then has to be mitigated downstream by threshold calibration and symbolic refinement.

A third limitation is that the symbolic layer, while valuable, introduces scalability trade-offs. The current implementation is entirely reasonable for hourly and moderate-scale analysis, but richer rule sets, higher-frequency data, or broader geographical coverage would increase grounding and reasoning complexity. Thus, the thesis demonstrates feasibility in a realistic but still bounded operating regime rather than proving unrestricted large-scale symbolic deployment.

A fourth limitation concerns the finance-aware evaluation. The financial utility analysis is an important step toward practical relevance, but it should still be interpreted as a \emph{simulated usefulness proxy} rather than as proof of a production-ready market strategy. The mapping from anomalies to utility is informative and methodologically valuable, but it remains a simplified backtesting layer rather than a full operational decision engine.

A fifth limitation is that some parts of the interpretability layer are currently strongest at the system level rather than at the level of fully explicit natural-language solver explanations. In particular, the report-level veto categories are partly reconstructed through repository logic and summary scripts rather than exposed directly as rich textual explanations from the ASP solver output. This still provides substantial interpretive value, but it also indicates a direction for further refinement.

These limitations do not weaken the contribution of the thesis. Rather, they define its scope precisely. The work should be understood as a strong demonstration of a lightweight neuro-symbolic anomaly-analysis framework, not as a final or exhaustive solution to all challenges in market-scale anomaly intelligence.

\subsection{Future Research}
\label{subsec:future_research}

The results of this thesis suggest several promising directions for future work.

A first and especially relevant direction is the integration of \emph{Large Language Models (LLMs)} into the neuro-symbolic pipeline. One natural extension would be to use LLMs as assistants for symbolic knowledge engineering. In such a setup, an LLM could support the drafting, translation, or refinement of candidate ASP rules from domain documentation, analyst notes, or high-level natural-language constraints. This would not remove the need for expert validation, but it could reduce the manual burden of rule creation and help scale the symbolic layer more efficiently. A second LLM-related direction would be to use language models not for rule generation itself, but for \emph{rule application support}, such as generating human-readable justifications, summarizing veto patterns, or producing analyst-facing explanations of why certain anomaly candidates were rejected \cite{openai2023gpt4,brown2020language}.

A second future direction is \emph{broader geographical generalization}. The current thesis focuses on Austria, which is appropriate for a master's-scale study and provides a coherent testbed with manageable complexity. However, the same methodological framework could be extended to broader European grids or multi-country market regions. Such an extension would be scientifically valuable because it would test whether the combination of lightweight learning and symbolic reasoning remains useful when cross-border coupling, heterogeneous weather regimes, and broader market interactions are introduced. It would also provide a much richer setting for evaluating multi-signal and multi-zone anomaly propagation.

A third direction is \emph{stronger forecasting backbones under the same deployment constraints}. Since the current results show that forecasting remains a challenging part of the pipeline, especially under volatile conditions, future work could investigate whether improved lightweight architectures or hybrid ensembles can preserve the compactness of dCeNN--ELM while improving predictive fit. This would be especially interesting if done without sacrificing edge feasibility or modularity.

A fourth direction is \emph{incremental or scalable symbolic reasoning}. For broader deployments, it would be useful to explore incremental ASP, window-based reasoning, or hierarchical symbolic filtering architectures. Such approaches could preserve the interpretability benefits of the current framework while reducing the computational burden associated with larger-scale or more time-sensitive analysis.

A fifth direction is \emph{richer downstream decision evaluation}. The finance-aware backtesting stage already moves the thesis beyond raw anomaly counts, but future work could extend this idea into more realistic decision-support settings. For example, anomaly signals could be integrated with probabilistic policies, operator dashboards, uncertainty intervals, or multi-step planning logic.

Finally, a particularly promising long-term direction is the development of \emph{interactive neuro-symbolic anomaly systems} in which statistical detection, symbolic logic, and language-based explanation are brought together in one analyst-facing workflow. In such a system, the forecasting backbone would generate anomaly candidates, ASP would apply hard plausibility constraints, and an LLM-like interface could communicate the resulting reasoning process in a way that is more accessible to non-technical users while still preserving formal traceability.

\subsection{Closing Statement}
\label{subsec:closing_statement}

In conclusion, this thesis has shown that lightweight neuro-symbolic anomaly analysis for weather-dependent electricity systems is both feasible and meaningful. The proposed framework does not claim universal superiority in every isolated metric. Instead, it demonstrates something more valuable: that a compact forecasting layer, when paired with explicit symbolic reasoning and practical downstream evaluation, can produce anomaly signals that are more interpretable, more plausibility-aware, and more operationally useful than purely statistical detection alone.

That is the central contribution of the work. The thesis establishes a bridge between learning and reasoning in a domain where both are needed: learning to capture complex multivariate behaviour under uncertainty, and reasoning to ensure that the resulting anomaly signals remain defensible in physical, market, and decision-oriented terms. In that sense, the research has achieved its main objective and provides a strong foundation for future neuro-symbolic anomaly-analysis research in electricity systems and beyond.